\section{Introducción}

La segregación residencial es un fenómeno emergente en el que numerosas decisiones locales pueden producir distribuciones espaciales globales marcadas. El modelo de Schelling muestra cómo preferencias individuales moderadas pueden generar segregación a escala colectiva \cite{schelling1971}.

Este proyecto adapta el modelo a una representación simplificada de Montevideo e incorpora factores sociales y económicos. Su escala y la necesidad de evaluar múltiples escenarios motivan el uso de computación de alto desempeño.

\subsection{Objetivos y contribuciones}

Las contribuciones previstas son una especificación reproducible del modelo, una implementación secuencial de referencia, una versión híbrida MPI/OpenMP y una evaluación de corrección, desempeño y escalabilidad. \pendiente{ajustar las contribuciones a lo efectivamente implementado.}

\subsection{Motivación computacional}

La evaluación repetida de vecindarios y destinos para cientos de miles de hogares, durante numerosos meses y escenarios, produce un costo suficiente para estudiar paralelismo de memoria distribuida y compartida. \pendiente{cuantificar el costo observado y evitar justificar HPC solamente por el tamaño nominal.}

\subsection{Organización del documento}

Las secciones siguientes describen el modelo y los datos, las implementaciones, la metodología de evaluación, los resultados y las conclusiones.
